\section{Simple period zeros and finite simultaneous choices}\label{sec:prelim}

We use polynomial one-forms $\omega=A\,dx+B\,dy$, with
$\deg\omega=\max(\deg A,\deg B)$. A period zero always refers to a
specified connected family of regular compact ovals, parametrized by the
Hamiltonian value. Distinct ovals can have the same value. We select
pairwise disjoint small closed annuli about the finitely many ovals that
will produce cycles.

\begin{lemma}\label{lem:jets}
Let $f_1,\ldots,f_M$ be linearly independent real analytic functions on a
connected open interval $I$. For every nonempty open subinterval $J\subset I$
there is a real linear combination with at least $M-1$ distinct simple
zeros in $J$.
\end{lemma}
\begin{proof}
For analytic functions linear independence implies that the Wronskian is
not identically zero. One way to see this is to perform elimination on
their convergent Taylor series at an interior point: independent analytic
germs admit a basis with strictly increasing orders of vanishing, whose
Wronskian has a nonzero leading term. The identity theorem transfers
independence to every open subinterval. Choose a point $h_0$ there where
the Wronskian is nonzero. The map from coefficients to the Taylor
coefficients through order $M-1$ at $h_0$ is invertible.

Fix distinct real numbers $t_1,\ldots,t_{M-1}$ and prescribe those Taylor
coefficients to be the coefficients of
$\prod_{j=1}^{M-1}(h-h_0-\varepsilon t_j)$. The resulting coefficient
vector stays bounded as $\varepsilon\to0$. For the corresponding
function $F_\varepsilon$,
\[
 \varepsilon^{-(M-1)}F_\varepsilon(h_0+\varepsilon t)
 =\prod_{j=1}^{M-1}(t-t_j)+O(\varepsilon)
\]
in $C^1$ on each fixed compact $t$-interval. Each simple root persists,
and its energy lies in the prescribed subinterval for sufficiently small
$\varepsilon>0$. For $M=1$ the assertion is empty.
\end{proof}

\begin{lemma}\label{lem:persistence}
For a finite family of regular compact Hamiltonian annuli and fixed
smooth one-forms, selected simple period zeros persist under a sufficiently
small smooth change of the Hamiltonian and forms. A single smallness
threshold can be used for the whole finite family.
\end{lemma}
\begin{proof}
On each compact annulus choose a transverse coordinate and smoothly
parametrize nearby level curves; regularity gives this parametrization
by the implicit function theorem. Differentiating the integral in the
energy parameter shows $C^1$ convergence of the period functions when
the Hamiltonian and forms converge in sufficiently high finite
differentiability norms on a neighborhood of the annulus.
For each simple zero choose a closed energy interval $J$ on which
$|M'|\ge a_J>0$ and the two endpoint values have opposite signs and
absolute value at least $b_J>0$. Taking the minimum of these finitely
many positive margins preserves both properties under a common
$C^1$ perturbation. The intermediate value theorem and monotonicity
give a simple zero in each selected interval. The annuli remain disjoint
under a sufficiently small perturbation.
\end{proof}

The common threshold in this lemma is taken only after the degree and
all the finite forms under consideration have been fixed. No threshold
uniform in the degree is used.

The following first-order criterion is standard; see, for example,
\cite[p.~443]{LiLiLlibreZhang2001}. We record its finite simultaneous form.

\begin{lemma}\label{lem:melnikov}
Suppose that a polynomial Hamiltonian $H$ and a polynomial one-form
$\omega=A\,dx+B\,dy$ have $N$ selected simple period zeros in pairwise
disjoint compact regular annuli. For all sufficiently small nonzero real
$\mu$, the single polynomial field
\begin{equation}\label{eq:singlefield}
 X_\mu=(H_y-\mu B,\,-H_x+\mu A)
\end{equation}
has at least $N$ distinct hyperbolic limit cycles and degree at most
$\max(\deg H-1,\deg\omega)$.
\end{lemma}
\begin{proof}
In the energy coordinate on a local transversal the return displacement
is $\mu M(h)+O(\mu^2)$ in $C^1$, where
$M(h)=\oint_{H=h}\omega$ is oriented by the Hamiltonian flow.
This expansion follows by smooth dependence of the flow and return time
and differentiation at $\mu=0$. Apply Lemma~\ref{lem:persistence} to the
displacement divided by $\mu$, using the minimum of the finite margins.
At each resulting fixed point the derivative of the return map is
$1+\mu M'(h)+O(\mu^2)$ and differs from $1$ for sufficiently small
nonzero $\mu$. It is positive and close to $1$, so the cycle is
hyperbolic. The degree follows directly from \eqref{eq:singlefield}.
\end{proof}
